% tex8web: tcolorbox's boxes, as near the PDF as CSS draws them: the frame
% (rules, round or sharp corners, borderlines, shadow), the colors, with
% darker ones for dark mode, the space around and inside, and the title (its
% font, a boxed title, a title that starts the text). tex4ht's own boxes keep
% only the colors and padding. What a skin draws freehand with TikZ
% (interior code, frame code, shadings) comes out as a plain frame. The
% static styles are in css/tex8web.css.
\makeatletter
\@ifpackageloaded{tcolorbox}{}{\makeatother\endinput}

% --- what tex4ht leaves out -------------------------------------------------

% the measures of a boxed title, which tex4ht never sets (styles use them:
% yshift=-\tcboxedtitleheight)
\providecommand\tcboxedtitleheight{0pt}
\providecommand\tcboxedtitlewidth{0pt}

% tex4ht closes a box before tcolorbox adds its after upper (a proof's
% \hfill$\square$), which then follows the box: add it first
\web@needs{tcb@endboxanddraw}
\let\web@tcb@endboxanddraw\tcb@endboxanddraw
\def\tcb@endboxanddraw{%
    \tcb@insert@after@part
    \let\tcb@insert@after@part\relax
    \web@tcb@endboxanddraw}
% and what follows an \hfill in an after upper or lower goes at the right:
% in macro #1, \web@hfillright before #2
\def\web@tcb@hfillafter#1#2{%
    \patchcmd#1{#2}{\web@hfillright#2}{}{\web@outdated{\string#1}}}
\web@tcb@hfillafter\tcb@insert@after@upper@unskip\kvtcb@after@upper
\web@tcb@hfillafter\tcb@insert@after@upper@nounskip\kvtcb@after@upper
\web@tcb@hfillafter\tcb@insert@after@lower@unskip\kvtcb@after@lower
\web@tcb@hfillafter\tcb@insert@after@lower@nounskip\kvtcb@after@lower
% #1, with what follows its first \hfill in a span css/tex8web.css puts at
% the right
\let\web@endhfill\relax
\newcommand\web@hfillright[1]{%
    \begingroup
        \def\hfill{\ifx\web@endhfill\relax
            \HCode{<span class="tex8web-hfill">}%
            \gdef\web@endhfill{\HCode{</span>}}%
        \fi}%
        #1%
        \web@endhfill
        \global\let\web@endhfill\relax
    \endgroup
}

% --- a box's style ----------------------------------------------------------

% (the options that leave nothing to read now, such as frame hidden, are
% noted as they're set, by tex8web-hooks.sty; a frame or interior is also
% hidden when the skin doesn't draw it, as with blanker)

% custom properties on the box's element, collected in \web@tcb@css (each
% helper appends one or more):
% - colors: a color in light mode, and in dark mode the same color at
%   another lightness (l, of oklch): frames and lines at least 0.62, so they
%   show; backgrounds dark as much as they were light; text light as much
%   as it was dark, or the page's text color for black
\def\web@tcb@add#1{\xdef\web@tcb@css{\web@tcb@css#1}}
\def\web@tcb@lightdark#1#2{light-dark(\##1, oklch(from \##1 #2 c h))}
\def\web@tcb@color#1#2#3{% name, color, dark lightness
    \web@hex{#2}\web@tcb@hex
    \web@tcb@add{--tcb-#1: \web@tcb@lightdark\web@tcb@hex{#3};}}
\def\web@tcb@frame#1#2{\web@tcb@color{#1}{#2}{max(l, 0.62)}}
\def\web@tcb@back#1#2{\web@tcb@color{#1}{#2}{min(l, 1.24 - l)}}
\def\web@tcb@black{000000}
\def\web@tcb@text#1#2{%
    \web@hex{#2}\web@tcb@hex
    \ifx\web@tcb@hex\web@tcb@black
        \web@tcb@add{--tcb-#1: var(--tex8web-text);}%
    \else
        \web@tcb@color{#1}{#2}{max(l, 1 - l)}%
    \fi}
\def\web@tcb@transparent#1{\web@tcb@add{--tcb-#1: transparent;}}
% - lengths, in ems (of a glue, its natural length): four of them, for the
%   sides (top, right, bottom, left) or the corners (from the top left);
%   the left side also on its own (-left)
\def\web@tcb@em#1{\web@em{\dimexpr\glueexpr#1\relax\relax}}
\def\web@tcb@sides#1#2#3#4#5{%
    \web@tcb@add{--tcb-#1: \web@tcb@em{#2} \web@tcb@em{#3} \web@tcb@em{#4} \web@tcb@em{#5};
        --tcb-#1-left: \web@tcb@em{#5};}}
\def\web@tcb@length#1#2{\web@tcb@add{--tcb-#1: \web@tcb@em{#2};}}
\def\web@tcb@sharp{sharp}
\def\web@tcb@corner#1{\ifx#1\web@tcb@sharp 0em\else\web@tcb@em\kvtcb@arc\fi}
% - the frame, the corners and the shadow of the box being drawn, under
%   names that start with #1
\def\web@tcb@frameof#1{%
    \ifx\tcb@frame@code\@empty\let\web@tcb@framehidden\@empty\fi
    \ifx\web@tcb@framehidden\relax
        \web@tcb@frame{#1frame}{tcbcolframe}%
    \else
        \web@tcb@transparent{#1frame}%
    \fi
    \web@tcb@length{#1rule-top}\kvtcb@top@rule@stand
    \web@tcb@length{#1rule-right}\kvtcb@right@rule
    \web@tcb@length{#1rule-bottom}\kvtcb@bottom@rule@stand
    \web@tcb@length{#1rule-left}\kvtcb@left@rule
    \web@tcb@add{--tcb-#1radius: \web@tcb@corner\tcb@corner@mode@NW\space
        \web@tcb@corner\tcb@corner@mode@NE\space
        \web@tcb@corner\tcb@corner@mode@SE\space
        \web@tcb@corner\tcb@corner@mode@SW;}%
    \ifcsname kvtcb@shadow\endcsname\ifx\kvtcb@shadow\@empty\else
        \web@tcb@add{--tcb-#1shadow: 0.2em 0.25em 0.4em rgb(0 0 0 / 0.3);}%
    \fi\fi
}
% - borderlines, which take the place of the frame on their side (the last
%   one there, if there are several)
\def\web@tcb@borderlines{%
    \ifcsname kvtcb@borderline\endcsname
        \begingroup
            \def\tcb@border##1##2##3{%
                \web@tcb@line{top}{##1}{##3}\web@tcb@line{right}{##1}{##3}%
                \web@tcb@line{bottom}{##1}{##3}\web@tcb@line{left}{##1}{##3}}%
            \def\tcb@bordernorth##1##2##3{\web@tcb@line{top}{##1}{##3}}%
            \def\tcb@bordereast##1##2##3{\web@tcb@line{right}{##1}{##3}}%
            \def\tcb@bordersouth##1##2##3{\web@tcb@line{bottom}{##1}{##3}}%
            \def\tcb@borderwest##1##2##3{\web@tcb@line{left}{##1}{##3}}%
            \kvtcb@borderline
        \endgroup
    \fi
}
% (a borderline's last argument is TikZ options, which are often just a
% color; other options get the frame's color)
\def\web@tcb@line#1#2#3{%
    \def\web@tcb@linecolor{tcbcolframe}%
    \in@{=}{#3}\ifin@\else\in@{,}{#3}\ifin@\else
        \in@{!}{#3}%
        \ifin@\def\web@tcb@linecolor{#3}\else\@ifundefinedcolor{#3}{}{\def\web@tcb@linecolor{#3}}\fi
    \fi\fi
    \expandafter\web@hex\expandafter{\web@tcb@linecolor}\web@tcb@hex
    \web@tcb@add{--tcb-line-#1: \web@tcb@em{#2} solid
        \web@tcb@lightdark\web@tcb@hex{max(l, 0.62)};}%
}

% the box's whole style, and its classes
\def\web@tcb@style{%
    \gdef\web@tcb@css{}%
    \web@tcb@frameof{}%
    \web@tcb@borderlines
    \ifx\tcb@interior@code\@empty\let\web@tcb@interiorhidden\@empty\fi
    \ifx\web@tcb@interiorhidden\relax
        \web@tcb@back{back}{tcbcolback}%
    \else
        \web@tcb@transparent{back}%
    \fi
    \web@tcb@text{text}{tcbcolupper}%
    \web@tcb@sides{padding}{\kvtcb@boxsep+\kvtcb@top}{\kvtcb@boxsep+\kvtcb@rightupper}%
        {\kvtcb@boxsep+\kvtcb@bottom}{\kvtcb@boxsep+\kvtcb@leftupper}%
    \ifx\web@tcb@before\relax\else\web@tcb@length{before}\web@tcb@before\fi
    \ifx\web@tcb@after\relax\else\web@tcb@length{after}\web@tcb@after\fi
    % the title: its text's color, and its background, the frame's color
    % unless it has its own; a boxed title is a box of its own, which
    % tcolorbox draws with the frame's color, the title's background and
    % the boxed title style
    \web@tcb@text{title-text}{tcbcoltitle}%
    \iftcb@titlefilled
        \web@tcb@back{title-back}{tcbcolbacktitle}%
    \else
        \web@tcb@back{title-back}{tcbcolframe}%
    \fi
    \web@tcb@sides{title-padding}{\kvtcb@boxsep+\kvtcb@toptitle}%
        {\kvtcb@boxsep+\kvtcb@righttitle}{\kvtcb@boxsep+\kvtcb@bottomtitle}%
        {\kvtcb@boxsep+\kvtcb@lefttitle}%
    \ifx\web@tcb@boxed\relax\else
        \begingroup
            \colorlet{web@tcb@frame}{tcbcolframe}%
            \colorlet{web@tcb@back}{tcbcolbacktitle}%
            \let\web@tcb@boxedstyle\kvtcb@style@boxedtitle
            \tcbset{reset,size=title,colframe=web@tcb@frame,colback=web@tcb@back}%
            \expandafter\tcbset\expandafter{\web@tcb@boxedstyle}%
            \web@tcb@frameof{title-}%
            \ifx\tcb@interior@code\@empty\let\web@tcb@interiorhidden\@empty\fi
            % (a boxed title that doesn't draw its interior has most likely
            % a frame code of its own that fills it)
            \ifx\web@tcb@interiorhidden\relax
                \web@tcb@back{title-back}{tcbcolback}%
            \else\ifx\web@tcb@framehidden\relax
                \web@tcb@back{title-back}{tcbcolframe}%
            \else
                \web@tcb@transparent{title-back}%
            \fi\fi
            \web@tcb@sides{title-padding}{\kvtcb@boxsep+\kvtcb@top}%
                {\kvtcb@boxsep+\kvtcb@rightupper}{\kvtcb@boxsep+\kvtcb@bottom}%
                {\kvtcb@boxsep+\kvtcb@leftupper}%
        \endgroup
    \fi
    \edef\web@tcb@classes{%
        \ifx\web@tcb@boxed\relax\else\space tex8web-tcb-boxed-\web@tcb@boxed\fi
        \ifx\web@tcb@inline\relax\else\space tex8web-tcb-inline\fi}%
    \web@Css{\#\csname :tcbcolid\endcsname\space{\web@tcb@css}}%
}

% whether the box steps a counter (\web@tcb@numbered empty), its own or
% one named by step, which tcolorbox does in the box's phantom; tex4ht never
% puts the phantom on the page, so the counter's anchor is lost. (A raster
% steps a counter of its own there for each box, which doesn't count)
\web@needs{tcb@set@@phantom}
\let\web@tcb@setphantom\tcb@set@@phantom
\ExplSyntaxOn
\str_new:N \l__web_phantom_str
\def\tcb@set@@phantom{%
    \str_set:NV \l__web_phantom_str \kvtcb@phantom
    \str_remove_all:Nn \l__web_phantom_str { \refstepcounter { tcbrasternum } }
    \str_if_in:NnTF \l__web_phantom_str { \refstepcounter }
        { \let\web@tcb@numbered\@empty } { \let\web@tcb@numbered\relax }
    \web@tcb@setphantom
    \web@tcb@keep}
\ExplSyntaxOff

% tcolorbox.4ht keeps a box's label, from its phantom to its drawing, after
% its contents, in one global macro, which a box inside it takes or
% replaces with its own. So the box keeps it, locally, until it's drawn
% (\Configure{tcolorbox}); and its number for the title (a theorem's
% \thetcbcounter), which tex4ht writes then too, after a box inside it has
% stepped the counter
\let\web@tcb@label\relax
\def\web@tcb@keep{%
    \ifcsname tcolorbox:label:key\endcsname
        \web@tcb@letcs\web@tcb@label{tcolorbox:label:key}%
        \web@tcb@letcs\web@tcb@currentlabel{tcolorbox:currentlabel}%
        \web@tcb@letcs\web@tcb@crefcurrentlabel{tcolorbox:cref@currentlabel}%
        \global\expandafter\let\csname tcolorbox:label:key\endcsname\undefined
    \else
        \let\web@tcb@label\relax
    \fi
    \ifdefined\thetcbcounter\protected@edef\thetcbcounter{\thetcbcounter}\fi}
\def\web@tcb@give{%
    \ifx\web@tcb@label\relax\else
        \global\expandafter\let\csname tcolorbox:label:key\endcsname\web@tcb@label
        \global\expandafter\let\csname tcolorbox:currentlabel\endcsname\web@tcb@currentlabel
        \global\expandafter\let\csname tcolorbox:cref@currentlabel\endcsname\web@tcb@crefcurrentlabel
    \fi}
% \let#1 to the macro named #2, which may be undefined (\csname would make
% it \relax)
\def\web@tcb@letcs#1#2{%
    \ifcsname#2\endcsname
        \expandafter\let\expandafter#1\csname#2\endcsname
    \else
        \let#1\undefined
    \fi}

% tex4ht's box, with that style; the title in its font (fonttitle), which
% tex4ht leaves out; and an anchor for a numbered box, where tcolorbox.4ht
% puts a label's, for page.lua's link. A \tcbox in its line (on line, or
% another way to have nothing before and after it) is in its line, as in
% the PDF: spans in the paragraph, which it doesn't end. (Its end comes
% after tcolorbox's group, so its start pushes its tag)
\def\web@tcb@hbox{hbox}
\def\web@tcb@span{span}
\ExplSyntaxOn
\seq_new:N \g__web_tcb_tags_seq
\cs_new_protected:Npn \web@tcb@pushtag { \seq_gpush:NV \g__web_tcb_tags_seq \web@tcb@tag }
\cs_new_protected:Npn \web@tcb@poptag { \seq_gpop:NN \g__web_tcb_tags_seq \web@tcb@tag }
\ExplSyntaxOff
\newif\ifweb@tcb@hbox
\Configure{tcolorbox}
    {\def\web@tcb@tag{div}%
     \let\web@tcb@env\@currenvir
     \web@tcb@hboxfalse
     \ifx\kvtcb@capture\web@tcb@hbox\ifx\kvtcb@beforebox\@empty
         \web@tcb@hboxtrue
         \let\web@tcb@tag\web@tcb@span
         \def\web@tcb@env{tcbox}%
     \fi\fi
     \ifweb@tcb@hbox\else\ifvmode\IgnorePar\fi\EndP\fi
     \web@tcb@pushtag
     \web@tcb@give
     \web@tcb@style
     \edef\kvtcb@before@title{\unexpanded\expandafter{\kvtcb@fonttitle}%
         \unexpanded\expandafter{\kvtcb@before@title}}%
     \HCode{<\web@tcb@tag\space class="tcolorbox \web@tcb@env\web@tcb@classes"
         id="\csname :tcbcolid\endcsname">}%
     \ifx\web@tcb@numbered\@empty
         \ifx\web@tcb@label\relax\AnchorLabel\fi
     \fi}
    {\ifx\kvtcb@title\@empty\else
         \HCode{\Hnewline<\web@tcb@tag\space class="tcolorbox-title">}\ifweb@tcb@hbox\else\par\fi
     \fi}
    {\ifx\kvtcb@title\@empty\else
         \ifweb@tcb@hbox\else\ifvmode\IgnorePar\fi\EndP\fi\HCode{</\web@tcb@tag>}%
     \fi
     \HCode{\Hnewline<\web@tcb@tag\space class="tcolorbox-content">}}
    {\web@tcb@poptag
     \ifx\web@tcb@tag\web@tcb@span
         \HCode{</span></span>}%
     \else
         \ifvmode\IgnorePar\fi\EndP\HCode{\Hnewline</div>\Hnewline</div>}\par\ShowPar
     \fi}
\Configure{tcolorlowerbox}
    {\ifvmode\IgnorePar\fi\EndP\HCode{\Hnewline<div class="lowerbox">}}
    {\ifvmode\IgnorePar\fi\EndP\HCode{\Hnewline</div>}}

\makeatother
